% Revision (docs/PAPER1_REVISION_PLAN.md): comparison models and evaluation added after the first submission.

\paragraph{Comparison models added in the revision}
\emph{Lift model}: the world model with the roll-forward switched off, so the
same encoder, readout, size and training steps, and the plan enters only
through the readout of each period. It is the transition ablation: any
difference between the two models comes from the carried state.
\emph{Chronos-Bolt}~\cite{chronos}: the base checkpoint, zero-shot, given each
series' history of units; we use the mean of its quantiles as the point
forecast. \emph{Temporal Fusion Transformer}~\cite{tft}: one model per data
set on Dominick's and Frat, fitted on the last 156 weeks with the plan as
known future inputs and a negative binomial likelihood, 1,500 steps; on M5 its
training cost is beyond the compute available, and it is not run there.
\emph{Equal-weight ensemble}: the mean of the world model, LightGBM and
TimesFM with the plan, fixed before any revision result, with no weights
fitted. Neural models are one fit (seed 1) in the rolling windows; the final
windows keep the three-seed results reported before.

\paragraph{Rolling-origin windows}
To test whether the final-window results hold, each data set gains windows
taken before the development windows, so none was used for any choice:
six 8-week windows on each Dominick's category, four on Frat, and four 28-day
windows on M5 (November 2015 to February 2016), each forecast from everything
earlier. Every model runs in every window with the settings of its final
configuration, unchanged. These windows and every analysis below were written
down before any of them ran.

\paragraph{Comparing two models}
For a window and a level, the difference in WRMSSE between two models gets a
95\% interval from 1,000 resamples of items: an item drawn twice counts twice
in every series and aggregate it feeds, and the revenue weights are recomputed,
so the interval reflects how much the comparison depends on which items are
in the data. Across windows, the mean difference gets a 95\% interval from
5,000 resamples of windows. A model is said to beat another only when that
interval excludes zero. For each data set and level we also report the 90\%
model confidence set of Hansen, Lunde and Nason~\cite{hansen2011}: the models
that cannot be told apart from the best, by the range statistic with a
bootstrap over windows.

\paragraph{Intervals}
Each model's interval forecasts come from its own errors in the earlier
windows of the same data set: errors scaled by each series' RMSSE scale are
pooled, and their 2.5\% to 97.5\% quantiles at each horizon are added to the
point forecast (rolling conformal calibration). The first rolling window has no
earlier window and is not scored. We report the scaled pinball loss averaged
over the quantiles, weighted by revenue as in the M5 uncertainty
competition~\cite{m5uncertainty}, and the share of actual sales inside the
50\%, 80\% and 95\% intervals. TimesFM, Chronos-Bolt and the Temporal Fusion
Transformer also give quantiles of their own, which are scored the same way.

\paragraph{Decisions}
Two simulations turn forecasts into costs. \emph{Replenishment}: for each
item-store-period, order the model's interval quantile at a critical ratio of
0.5, 0.8 or 0.95; each unit over costs 1 and each unit short costs
$r/(1-r)$, summed over the window. \emph{Category planning}: the absolute error
of each category-store total over the window as a share of its sales, weighted
by dollar sales. Both are reported against seasonal naive. They are
simulations under stated costs, not business results.

\paragraph{Reconciliation}
Base forecasts at every node of the hierarchy come from AutoETS fitted on the
history before each window (seasonal at the aggregate nodes, non-seasonal at
item by store). They are reconciled bottom-up, top-down by historical
proportions, and by the MinT family~\cite{mint}: ordinary least
squares, weighted least squares with structural weights and with in-sample
error variances, and MinT with a shrunk covariance where the hierarchy has at
most 6,000 nodes. Each model's own item-by-store forecasts are also combined
with ETS above them by weighted least squares, to see whether reconciliation
improves the upper levels without losing the bottom.

\paragraph{Mechanism tests}
Three tests on the saved world models, with no refitting. \emph{Timing}: the
plan of the forecast periods is shuffled within each series, keeping its total
promotion exposure; a model that uses the timing of the plan loses accuracy.
\emph{Persistence}: for series with nothing planned in the window, a
promotion and a 20\% price cut are added in the first period, and the change
in the forecast is read in every later period; the lift model has no carried
state, so its later change is zero by construction. \emph{State}: a ridge
read-out from the rolled state to sales in the periods after the window,
against the same read-out from the state before the window.
